\chapter{Production de la machine, vérification humaine}\label{ch:review}

Dans un projet où un assistant produit beaucoup, la question la plus importante devient :
qu'est-ce qui a été vérifié, et par qui ? Si l'on ne peut plus y répondre, tout le travail
perd sa valeur, y compris la part qui était juste.

\section{Un en-tête sur chaque fichier}

\begin{regle}
Chaque fichier produit par un assistant dit, dans son en-tête, quel modèle l'a produit,
quand, à partir de quoi, et ce qui en a été vérifié.
\end{regle}

\begin{terminal}
% Pass: <modèle>, 2026-10-31 — first pass,
%   unchecked against the pages by a human.
\end{terminal}

Quand une personne vérifie un point, l'en-tête l'enregistre : la ligne corrigée, la date,
l'ancienne lecture. L'histoire du dépôt garde le reste.

\section{Marquer le doute au lieu de le résoudre}

Dans une transcription, l'illisible se marque (\verb|\ill{}|), et la lecture incertaine
aussi (\verb|\uncertain{...}|). Ces marques ne sont pas des défauts : elles
disent exactement où une personne doit regarder. Une transcription sans doute marqué n'est
pas plus fiable, elle cache ses doutes.

\section{Une page de relecture}

Pour qu'une personne vérifie vite, on lui apporte les questions au lieu de lui demander de
tout relire. Une page web locale montre, pour chaque doute, l'image de la page à côté de la
lecture proposée, avec une question fermée et trois boutons : oui, non, incertain. Les
réponses partent ensuite en un bloc vers l'assistant, qui les applique et ferme les
questions.

\begin{exemple}
Dans une séance de relecture, la plupart des réponses confirment une lecture ; les autres corrigent un mot, ou signalent que
le contexte autour du mot était faux. L'assistant applique les corrections, retire les
questions réglées et repose plus précisément celles qui restent.
\end{exemple}

\begin{tipblock}
Une question bien posée se répond en quelques secondes : un mot, une case, une page. Une
question qui demande de relire un paragraphe entier se découpe en plusieurs.
\end{tipblock}

\section{Ne pas confondre les statuts}

Une lecture vérifiée par la machine (une compilation, un calcul) et une lecture vérifiée
par une personne (sur le document original) ne sont pas la même chose ; on les distingue
dans les fichiers comme dans les articles. Un article sérieux sur un tel projet dit combien
de points ont été relus par une personne, sur combien.

\section{Ce que d'autres mathématiciens pratiquent}

Le blog collectif \emph{Proofs and Prompts} (\url{https://proofsandprompts.com}) rassemble
depuis l'automne 2026 des témoignages de mathématiciens sur leur usage des modèles de
langage. Plusieurs pratiques qu'on y lit rejoignent celles de ce chapitre, et méritent
d'être reprises.

\begin{itemize}
  \item \textbf{Relire avant de soumettre.} Un billet anonyme (« Use LLMs to check your
        work ») tient pour un devoir professionnel de faire relire un article par un
        modèle avant de l'envoyer : coquilles, trous dans les démonstrations, attributions.
        Les rapporteurs n'ont pas à perdre leur temps sur ce qu'une lecture rapide aurait
        trouvé.
  \item \textbf{Traiter chaque signalement comme une affirmation à vérifier.} Alonso
        Castillo-Ramirez a fait relire quelques dizaines d'articles publiés et vérifié chaque
        objection à la main ; cinq étaient fondées, dont une erreur dans un de ses propres
        articles. Vérifier un signalement prend des minutes, certifier un article des
        semaines. Il propose un registre public des erreurs signalées et de leurs
        corrections.
  \item \textbf{Distinguer « correct » et « démontré ».} Stepan Nesterov propose de dire
        qu'un énoncé est \emph{correct} quand une preuve formelle existe, et
        \emph{démontré} quand une démonstration en langue naturelle a été comprise et
        acceptée par des experts. C'est la distinction de ce chapitre entre ce que la
        machine a vérifié et ce qu'une personne a relu.
  \item \textbf{Reconstruire une preuve opaque.} Timothy Chow conseille, devant une
        démonstration produite par un modèle, de supposer qu'elle combine des idées
        connues, et de la reconstruire à partir d'elles et de leur littérature plutôt que de
        l'avaler d'un bloc.
  \item \textbf{Faire tourner la boucle des contre-exemples.} Yushang Cheng décrit une
        boucle : retirer une hypothèse, chercher un contre-exemple, garder le
        contre-exemple comme nouvel objet, écrire les chaînes de dépendances. C'est la
        boucle du chapitre sur \LaTeX{} et Lean, vue du côté de l'exploration.
  \item \textbf{Déclarer précisément.} Maximilian Kaipel propose une taxonomie des tâches
        confiées à l'assistant et un paquet \LaTeX{}, \texttt{mathaidisclose}, qui produit la
        déclaration ; Hendrik De Bie, après une enquête sur des dépôts arXiv, recommande de
        dire quel modèle a servi, pour quelle section, et à quel titre (exploration,
        démonstration, relecture). András Juhász plaide pour que les universités hébergent
        des modèles ouverts, pour la confidentialité et l'indépendance.
\end{itemize}

\section{Répondre aux inquiétudes}

Le même blog publie des textes inquiets, souvent justes. Les règles de ce livre n'y
répondent pas toutes, mais elles en désamorcent plusieurs.

\begin{itemize}
  \item \textbf{L'erreur et l'excès de confiance.} Un texte produit par un modèle convainc, et
        un commentateur rapporte des taux d'hallucination élevés dans ses propres agents.
        Réponse : rien n'est tenu pour vrai sans une vérification nommée (la compilation
        Lean, un calcul, une personne sur l'original), et chaque fichier dit laquelle. Un
        signalement de l'assistant est une question, jamais un verdict.
  \item \textbf{L'opacité.} Une preuve de plusieurs millions de lignes, ou une démonstration sans
        explication, n'apprend rien. Réponse : la preuve formelle ne remplace pas la
        rédaction ; elle la garantit, et la rédaction en \LaTeX{} explique ce qu'elle
        démontre, étape par étape, en renvoyant au fichier.
  \item \textbf{La charge des rapporteurs.} Les relecteurs risquent de laisser passer ce qu'ils
        croient « déjà vérifié ». Réponse : distinguer dans le texte ce qui a été vérifié
        par la machine de ce qui a été lu par une personne, et dire sur combien de points
        porte la relecture humaine. Le rapporteur sait alors où regarder.
  \item \textbf{Le crédit et l'attribution.} Un modèle peut reprendre un lemme sans citer son
        auteur. Réponse : aucune affirmation de nouveauté, une recherche nommée pour chaque
        énoncé, et un statut « trouvé » qu'on garde quand on retrouve l'énoncé ailleurs
        (chapitre~\ref{ch:litterature}).
  \item \textbf{La formation.} La lutte avec un exercice forme l'intuition, et plusieurs auteurs
        craignent qu'une génération s'en passe. Réponse partielle seulement : ce livre
        s'adresse à qui a déjà cette formation, et la boucle de relecture oblige à
        comprendre chaque énoncé qu'on signe. Pour les étudiants, la question reste
        ouverte, et je n'ai pas de réponse.
  \item \textbf{La dépendance aux entreprises.} Les outils appartiennent à des sociétés privées.
        Réponse : tout ce qui compte vit dans le dépôt, en texte brut, et se relit sans
        l'assistant ; un autre outil, ou un modèle ouvert hébergé par une université,
        pourrait reprendre le travail là où il est.
\end{itemize}
